\section{USAMO 2019}
        \begin{enumerate}
        	\item (\textit{USAMO 2019}) \textbf{Note:} For any geometry problem whose statement begins with an asterisk $(*)$, the first page of the solution must be a large, in-scale, clearly labeled diagram. Failure to meet this requirement will result in an automatic 1-point deduction.


 Let $\mathbb{N}$ be the set of positive integers. A function $f:\mathbb{N}\to\mathbb{N}$ satisfies the equation 
 \[\underbrace{f(f(\ldots f}_{f(n)\text{ times}}(n)\ldots))=\frac{n^2}{f(f(n))}\]for all positive integers $n$. Given this information, determine all possible values of $f(1000)$.


 



\textbf{Solution}

Let $f^r(x)$ denote the result when $f$ is applied to $f^{r-1}(x)$, where $f^1(x)=f(x)$.$\hfill \break \hfill \break$If $f(p)=f(q)$, then $f^2(p)=f^2(q)$ and $f^{f(p)}(p)=f^{f(q)}(q)$

 $\implies p^2=f^2(p)\cdot f^{f(p)}(p)=f^2(q)\cdot f^{f(q)}(q)=q^2$

 $\implies p=\pm q$

 $\implies p=q$ since $p,q>0$.

 Therefore, $f$ is injective. It follows that $f^r$ is also injective.

 
Lemma 1: If $f^r(b)=a$ and $f(a)=a$, then $b=a$.

 
Proof:

 $f^r(b)=a=f^r(a)$ which implies $b=a$ by injectivity of $f^r$.

 
Lemma 2: If $f^2(m)=f^{f(m)}(m)=m$, and $m$ is odd, then $f(m)=m$.

 
Proof:

 Let $f(m)=k$. Since $f^2(m)=m$, $f(k)=m$. So, $f^2(k)=k$. $\newline f^2(k)\cdot f^{f(k)}(k)=k^2$.

 Since $k\neq0$, $f^{f(k)}(k)=k$

 $\implies f^m(k)=k$

 $\implies f^{gcd(m, 2)}(k)=k$

 $\implies f(k)=k$

 This proves Lemma 2.

 
I claim that $f(m)=m$ for all odd $m$.

 Otherwise, let $m$ be the least counterexample.

 Since $f^2(m)\cdot f^{f(m)}(m)=m^2$, either

 $(1) f^2(m)=k<m$, contradicted by Lemma 1 since $k$ is odd and $f^2(k)=k$.

 $(2) f^{f(m)}(m)=k<m$, also contradicted by Lemma 1 by similar logic.

 $(3) f^2(m)=m$ and $f^{f(m)}(m)=m$, which implies that $f(m)=m$ by Lemma 2.This proves the claim.

 
By injectivity, $f(1000)$ is not odd.I will prove that $f(1000)$ can be any even number, $x$. Let $f(1000)=x, f(x)=1000$, and $f(k)=k$ for all other $k$. If $n$ is equal to neither $1000$ nor $x$, then $f^2(n)\cdot f^{f(n)}(n)=n\cdot n=n^2$. This satisfies the given property.

 If $n$ is equal to $1000$ or $x$, then $f^2(n)\cdot f^{f(n)}(n)=n\cdot n=n^2$ since $f(n)$ is even and $f^2(n)=n$. This satisfies the given property.

 


 


 


	\item (\textit{USAMO 2019}) \textbf{Note:} For any geometry problem whose statement begins with an asterisk $(*)$, the first page of the solution must be a large, in-scale, clearly labeled diagram. Failure to meet this requirement will result in an automatic 1-point deduction.


 Let $ABCD$ be a cyclic quadrilateral satisfying $AD^2 + BC^2 = AB^2$. The diagonals of $ABCD$ intersect at $E$. Let $P$ be a point on side $\overline{AB}$ satisfying $\angle APD = \angle BPC$. Show that line $PE$ bisects $\overline{CD}$.


 



\textbf{Solution}

Let $PE \cap DC = M$. Also, let $N$ be the midpoint of $AB$.Note that only one point $P$ satisfies the given angle condition. With this in mind, construct $P'$ with the following properties:

 (1) $AP' \cdot AB = AD^2$

 (2) $BP' \cdot AB = CD^2$

 Claim: $P = P'$

 Proof:The conditions imply the similarities $ADP \sim ABD$ and $BCP \sim BAC$ whence $\measuredangle APD = \measuredangle BDA = \measuredangle BCA = \measuredangle CPB$ as desired. $\square$

 Claim: $PE$ is a symmedian in $AEB$

 Proof:We have 
 \begin{align*}  AP \cdot AB = AD^2 \iff AB^2 \cdot AP &= AD^2 \cdot AB \\ \iff \left( \frac{AB}{AD} \right)^2 &= \frac{AB}{AP} \\ \iff \left( \frac{AB}{AD} \right)^2 - 1 &= \frac{AB}{AP} - 1 \\ \iff \frac{AB^2 - AD^2}{AD^2} &= \frac{BP}{AP} \\ \iff \left(\frac{BC}{AD} \right)^2 &= \left(\frac{BE}{AE} \right)^2 = \frac{BP}{AP}  \end{align*}
as desired. $\square$

 Since $P$ is the isogonal conjugate of $N$, $\measuredangle PEA = \measuredangle MEC = \measuredangle BEN$. However $\measuredangle MEC = \measuredangle BEN$ implies that $M$ is the midpoint of $CD$ from similar triangles, so we are done. $\square$

 



\textbf{Solution 2}


                    \begin{figure}[H]
                        \centering
                        \includegraphics[width=0.5\textwidth]{Solution/USAMO_I_2019_2_Solution_Figure_1.png}
                    \end{figure}
                    
                    \begin{figure}[H]
                        \centering
                        \includegraphics[width=0.5\textwidth]{Solution/USAMO_I_2019_2_Solution_Figure_2.png}
                    \end{figure}
                    Let $\omega$ be the circle centered at $A$ with radius $AD.$

 Let $\Omega$ be the circle centered at $B$ with radius $BC.$

 We denote $I_\omega$ and $I_\Omega$ inversion with respect to $\omega$ and $\Omega,$ respectively.
 \[B'= I_\omega (B), C'= I_\omega (C), D = I_\omega (D) \implies\]
 \[AB' \cdot AB = AD^2,  \angle ACB = \angle AB'C'.\]
 \[A'= I_\Omega (A), D'= I_\Omega (D), C = I_\Omega (C) \implies\]
 \[BA' \cdot AB = BC^2,  \angle BDA = \angle BA'D'.\]Let $\theta$ be the circle $ABCD.$

 $I_\omega (\theta) = B'C'D,$ straight line, therefore 
 \[\angle AB'C' = \angle AB'D' = \angle ACB.\] $I_\Omega (\theta) = A'D'C,$ straight line, therefore
 \[\angle BA'D' = \angle BA'C = \angle BDA.\]$ABCD$ is cyclic $\implies \angle BA'C = \angle AB'D.$
 \[AB' + BA' = \frac {AD^2 + BC^2 }{AB} = AB \implies\] points $A'$ and $B'$ are coincide.

 Denote $A' = B' = Q \in AB.$

 Suppose, we move point $Q$ from $A$ to $B.$ Then $\angle AQD$ decreases monotonically, $\angle BQC$ increases monotonically. So, there is only one point where 
 \[\angle AQD = \angle BQC \implies P = Q.\]

 
 \[B = I_\omega (P), D' = I_\omega (D'), C' = I_\omega (C), A = I_\omega (\infty) \implies\]$\hspace{19mm} I_\omega (CD'P) = AC'D'B$ is cyclic.
 \[\angle ACD = \angle ABD = \angle CC'D \implies C' D' || CD \implies\]$\hspace{19mm} C'D'CD$ is  trapezoid.

 It is known that the intersection of the diagonals, intersection point of the lines containing the lateral sides of the trapezoid and the midpoints of two parallel sides are collinear.

 \textbf{vladimir.shelomovskii@gmail.com, vvsss}

 


	\item (\textit{USAMO 2019}) \textbf{Note:} For any geometry problem whose statement begins with an asterisk $(*)$, the first page of the solution must be a large, in-scale, clearly labeled diagram. Failure to meet this requirement will result in an automatic 1-point deduction.


 Let $K$ be the set of all positive integers that do not contain the digit $7$ in their base-$10$ representation. Find all polynomials $f$ with nonnegative integer coefficients such that $f(n)\in K$ whenever $n\in K$.


 



 



\textbf{Solution}

I claim the only such polynomials are of the form $f(n)=k$ where $k\in K$, or $f(n)=an+b$ where $a$ is a power of 10, $b\in K$, and $b<a$. Obviously, these polynomials satisfy the conditions. We now prove that no other polynomial works. That is, we prove that for any other polynomial $f$ with nonnegative coefficients, there is some $n\in K$ such that $f(n)\notin K$.

 We first prove the result for monomials, polynomials in which only one coefficient is nonzero. This is obvious for constant polynomials $f(n)=k\notin K$. The next simplest case is $f(n)=an$ with $a$ not a power of 10, and hence $\lg a$ is irrational. By Dirichlet's approximation theorem, the set of multiples of $\lg a$ is dense $\bmod\ 1$, and thus contains an element with fractional part in the interval $[\lg 7,\lg 8)$. In other words, there is a power of $a$ whose decimal expansion starts with a 7. Let $a^x$ be the smallest power of $a$ that is not in $K$. Obviously, $x>0$, so letting $n=a^{x-1}$ completes the proof of this part.

 
We have now proven that for any $a$ that is not a power of 10, there is some $n\in K$ such that $an\notin K$. We proceed to the case where $f(n)=an^x$ for $x>1$. This splits into 2 cases. If $ax$ is not a power of 10, then pick $m\in K$ such that $axm\notin K$. For any $y$, we have
 \[f(10^y+m)=a10^{yx}+axm*10^{y(x-1)}+...+am^x\]If we choose $y$ to be large enough, then the terms in the expression above will not interfere with each other, and the resulting number will contain a 7 in the decimal expansion, and thus not be in $K$.

 On the other hand, if $ax$ is a power of 10, then $a$ and $x$ are both powers of 10, and $x\ge10$. Obviously, $\frac12ax(x-1)$ is not a power of 10. By the previous step, which establishes the result for $x=2$, we can pick $m\in K$ such that $\frac12ax(x-1)m^2\notin K$. Then, for any $y$,
 \[f(10^y+m)=a10^{yx}+axm*10^{y(x-1)}+\frac12ax(x-1)m^2*10^{y(x-2)}+...+am^x\]Similarly, picking a sufficient large $y$ settles this case.

 
Now, we extend our proof to general polynomials. If a polynomial $f(n)=a_0+a_1n+a_2n^2+...+a_xn^x$ satisfies the conditions of the problem, then for any $m,y>0$:
 \[f(m*10^y)=a_xm^x*10^{yx}+...+a_1m*10^y+a_0\]Similarly, choosing $y$ to be sufficiently large results in the terms not interfering with each other. If $f$ contains any monomials that do not satisfy the conditions of the problem, then picking suitable $m$ and sufficiently large $y$ causes $f(m*10^y)$ to not be in $K$. Thus, $f$ is a linear polynomial of the form $ax+b$ where $a$ is 0 or a power of 10, and $b\in K$. It suffices to rule out those polynomials where $a>0$ and $b>a$. If this is the case, then since the digit of $b$ corresponding to $a$ is not 7, there must be a single-digit number $n$ such that the digit of $f(n)$ corresponding to $a$ is 7. Therefore, we are done. 

 -wzs26843545602

 


	\item (\textit{USAMO 2019}) \textbf{Note:} For any geometry problem whose statement begins with an asterisk $(*)$, the first page of the solution must be a large, in-scale, clearly labeled diagram. Failure to meet this requirement will result in an automatic 1-point deduction.


 Let $n$ be a nonnegative integer. Determine the number of ways that one can choose $(n+1)^2$ sets $S_{i,j}\subseteq\{1,2,\ldots,2n\}$, for integers $i,j$ with $0\leq i,j\leq n$, such that:
for all $0\leq i,j\leq n$, the set $S_{i,j}$ has $i+j$ elements; and
$S_{i,j}\subseteq S_{k,l}$ whenever $0\leq i\leq k\leq n$ and $0\leq j\leq l\leq n$.


 



\textbf{Solution 1}

Note that there are $(2n)!$ ways to choose $S_{1, 0}, S_{2, 0}... S_{n, 0}, S_{n, 1}, S_{n, 2}... S_{n, n}$, because there are $2n$ ways to choose which number $S_{1, 0}$ is, $2n-1$ ways to choose which number to append to make $S_{2, 0}$, $2n-2$ ways to choose which number to append to make $S_{3, 0}$, etc. After that, note that $S_{n-1, 1}$ contains the $n-1$ in $S_{n-1, 0}$ and 1 other element chosen from the 2 elements in $S_{n, 1}$ not in $S_{n-1, 0}$ so there are 2 ways for $S_{n-1, 1}$. By the same logic there are 2 ways for $S_{n-1, 2}$ as well so $2^n$ total ways for all $S_{n-1, j}$, so doing the same thing $n-1$ more times yields a final answer of $(2n)!\cdot 2^{\left(n^2\right)}$.

 -Stormersyle

 



\textbf{Solution 2}

There are $\frac{(2n)!}{2^n}$ ways to choose $S_{0,0}, S_{1,1} .... S_{n,n}$. Since, there are $\binom{2n}{2}$ ways to choose $S_{1,1}$, and after that, to generate $S_{2,2}$, you take $S_{1,1}$ and add 2 new elements, getting you $\binom{2n-2}{2}$ ways to generate $S_{2,2}$. And you can keep going down the line, and you get that there are $\frac{(2n)!}{2^n}$ ways to pick $S_{0,0}, S_{1,1} ... S_{n,n}$ Then we can fill out the rest of the gird. First, let’s prove a lemma.

 \textbf{Lemma} Claim: If we know what $S_{a,b}$ is and what $S_{a+1,b+1}$ is, then there are 2 choices for both $S_{a,b+1}$ and $S_{a+1,b}$.

 Proof: Note $a \leq a+1$ and $b \leq b+1$, so $S_{a,b}\subseteq S_{a+1,b+1}$. Let $A$ be a set that contains all the elements in $S_{a+1,b+1}$ that are not in $S_{a,b}$. $A = S_{a+1,b+1} - S_{a,b}$. We know $S_{a+1,b+1}$ contains total $a+b+2$ elements. And $S_{a,b}$ contains total $a+b$ elements. That means $A$ contains only 2 elements since $(a+b+2)-(a+b) = 2$. Let’s call these 2 elements $m, n$. $A = \{m, n\}$. $S_{a,b+1}$ contains 1 elements more than $S_{a,b}$ and 1 elements less than $S_{a+1,b+1}$. That 1 elements has to select from $A$. It’s easy to see $S_{a,b+1} = S_{a,b}\cup \{m\}$ or $S_{a,b+1} = S_{a,b}\cup \{n\}$, so there are 2 choice for $S_{a,b+1}$. Same thing applies to $S_{a+1,b}$.

 \textbf{Filling in the rest of the grid} We used our proved lemma, and we can fill in $S_{0,1}, S_{1,2} ... S_{n-1,n}$ then we can fill in the next diagonal, until all $S_{i,j}$ are filled, where $i \leq j$. But, we haven’t finished everything! Fortunately, filling out the rest of the diagonals in a similar fashion is pretty simple. And, it’s easy to see that we have made $n(n+1)$ decisions, each with 2 choices, when filling out the rest of the grid, so there are $2^{n(n+1)}$ ways to finish off.

 \textbf{Finishing off} To finish off, we have $\frac{(2n)!}{2^n} \cdot 2^{n(n+1)}$ ways to fill in the grid, which gets us $\boxed{(2n)! \cdot 2^{n^2}}$

 -Alexlikemath

 



\textbf{Solution 3}

Let $C_{j}$ represent the set of sets of the form $S_{ij}$ for $1 \le i \le n$, $a_{ij}$ denote $S_{(i+1)j} \backslash S_{ij}$, and $b_{ij} = S_{i(j+1)} \backslash S_{ij}$. Begin by considering $C_0$ and $S_{00} = \emptyset$. Then  given $S_{i0}$ we can create $S_{(i+1)0}$ by adding one element ($a_{i0}$). Using this, the number of ways to form the sequence of $S_{00}, S_{10}, \dots, S_{n0}$ are $(2n)(2n-1) \cdots (n+1)$ where we successively add one of the remaining elements of $[2n]$ to get consecutive terms in the sequence.

 Now consider when we are given $C_{j}$ and we need to find $C_{j+1}$. So far, there have been $n + j$ chosen distinct elements (via $S_{nj}$). After finding $C_{j+1}$ we will have $n + j + 1$ distinct elements and so in this process we only add one unique element to sets among $C_{j+1}$. There are $2n - (n+j) = n-j$ ways to chose such a new element called $x$.

 Now notice that $b_{0j}, a_{0(j+1)}, a_{1(j+1)}, \dots, a_{(n-1)(j+1)}$ is a permutation of $x,  a_{0j}, \dots, a_{(n-1)j}$ by noting $b_{nj} = x$ and,
 \[S_{0j} + b_{0j} +  a_{0(j+1)} +  a_{1(j+1)} +  \cdots +  a_{(n-1)(j+1)} = S_{n(j+1)} = S_{0j} +  a_{0j}+ a_{1j} +  \cdots +  a_{(n-1)j} + b_{nj}\]Furthermore,
 \begin{align*} S_{0j} + b_{0j} + a_{0(j+1)} &= S_{1(j+1)} =  S_{0j} +  a_{0j} +  b_{1j} \\ S_{1j} + b_{1j} + a_{1(j+1)} &= S_{2(j+1)} = S_{1j} +  a_{1j} +  b_{2j} \\ &\cdots \\ S_{(n-1)j} + b_{(n-1)j} + a_{(n-1)(j+1)} &= S_{n(j+1)} =  S_{(n-1)j} +  a_{(n-1)j} +  b_{nj} . \end{align*}
Therefore $b_{ij}, a_{i(j+1)}$ is a permutation of $a_{ij}, b_{(i+1)j}$ for $i < n$. Now let $k$ be the first $i$ such that $b_{ij} = x$. By definition, $b_{(k+1)j}, \dots = x$. Then the number of ways to order $x,  a_{0j}, \dots, a_{(n-1)j}$ is $2^{k-1}$ as there are 2 permutations for each pair before $\{b_{ij}, a_{i(j+1)}\} = \{a_{ij}, b_{(i+1)j}\}$ and each pair after is determined by $b_{ij}= b_{(i+1)j} = k, a_{i(j+1)} = a_{ij}$. For $k  = 0$, the permutation is completely determined so there is one way.

 Overall the number of ways to add the $j+1$'th row is, 
 \[(n-j)(1 + \sum_{k = 1}^n 2^{k-1}) = 2^n(n-j).\]

 In total, there are $(2n)(2n-1)\cdots (n+1)$ ways to find $C_0$ and for each $C_{j}$ there are $2^n(n-j+1)$ ways for $1 \le j \le n$. So the answer is, 
 \[\frac{2n!}{n!}(2^nn)(2^{n}(n-1)) \cdots (2^{n}1) = \boxed{2^{n^2} \cdot (2n)!}.\]

 


	\item (\textit{USAMO 2019}) \textbf{Note:} For any geometry problem whose statement begins with an asterisk $(*)$, the first page of the solution must be a large, in-scale, clearly labeled diagram. Failure to meet this requirement will result in an automatic 1-point deduction.


 Two rational numbers $\frac{m}{n}$ and $\frac{n}{m}$ are written on a blackboard, where $m$ and $n$ are relatively prime positive integers. At any point, Evan may pick two of the numbers $x$ and $y$ written on the board and write either their arithmetic mean $\frac{x+y}{2}$ or their harmonic mean $\frac{2xy}{x+y}$ on the board as well. Find all pairs $(m,n)$ such that Evan can write $1$ on the board in finitely many steps.


 



\textbf{Solution}

We claim that all odd $m, n$ work if $m+n$ is a positive power of 2.

 Proof:We first prove that $m+n=2^k$ works. By weighted averages we have that $\frac{n(\frac{m}{n})+(2^k-n)\frac{n}{m}}{2^k}=\frac{m+n}{2^k}=1$ can be written, so the solution set does indeed work. We will now prove these are the only solutions.

 Assume that $m+n\ne 2^k$, so then $m+n\equiv 0\pmod{p}$ for some odd prime $p$. Then $m\equiv -n\pmod{p}$, so $\frac{m}{n}\equiv \frac{n}{m}\equiv -1\pmod{p}$. We see that the arithmetic mean is $\frac{-1+(-1)}{2}\equiv -1\pmod{p}$ and the harmonic mean is $\frac{2(-1)(-1)}{-1+(-1)}\equiv -1\pmod{p}$, so if 1 can be written then $1\equiv -1\pmod{p}$ and $2\equiv 0\pmod{p}$ which is obviously impossible, and we are done.

 -Stormersyle   

 


	\item (\textit{USAMO 2019}) \textbf{Note:} For any geometry problem whose statement begins with an asterisk $(*)$, the first page of the solution must be a large, in-scale, clearly labeled diagram. Failure to meet this requirement will result in an automatic 1-point deduction.


 Find all polynomials $P$ with real coefficients such that 
 \[\frac{P(x)}{yz}+\frac{P(y)}{zx}+\frac{P(z)}{xy}=P(x-y)+P(y-z)+P(z-x)\]holds for all nonzero real numbers $x,y,z$ satisfying $2xyz=x+y+z$.


 



\textbf{Solution}

If $P(x)=c$ for a constant $c,$ then $\dfrac{c(x+y+z)}{xyz}=3c$. We have $2c=3c.$ Therefore $c=0.$

 Now consider the case of non-constant polynomials. First we have 
 \[xP(x)+yP(y)+zP(z)=xyz(P(x-y)+P(y-z)+P(z-x))\] for all nonzero real numbers $x,y,z$ satisfying $2xyz=x+y+z$. Both sides of the equality are polynomials (of $x,y,z$). They have the same values on the 2-dimensional surface $2xyz=x+y+z$, except for some 1-dimensional curves in it. By continuity, the equality holds for all points on the surface, including those with $z=0.$ Let $z=0,$ we have $y=-x$ and $x(P(x)-P(-x))=0.$ Therefore $P$ is an even function.   

 (Here is a sketch of an elementary proof. Let $z=\dfrac{x+y}{2xy-1}.$ We have
 \[xP(x)+yP(y)+\dfrac{x+y}{2xy-1}P(\dfrac{x+y}{2xy-1})=xy\dfrac{x+y}{2xy-1}(P(x-y)+P(y-\dfrac{x+y}{2xy-1})+P(\dfrac{x+y}{2xy-1}-x)).\]This is an equality of rational expressions. By multiplying $(2xy-1)^N$ on both sides for a sufficiently large $N$, they become polynomials, say $A(x,y)=B(x,y)$ for all real $x, y$ with $x\ne 0, y\ne 0, x+y\ne 0$ and $2xy-1\ne 0.$ For a fixed $x,$ we have two polynomials (of $y$) having same values for infinitely many $y$. They must be identical. Let $y=0,$ we have $x^{N+1}(P(x)-P(-x))=0.$ )

 Notice that if $P(x)$ is a solution,  then is $cP(x)$ for any constant $c.$ For simplicity, we assume the leading coefficient of $P$ is $1$: 
 \[P(x)=x^n+a_{n-2}x^{n-2}+\cdots +a_2x^2+a_0,\] where $n$ is a positive even number.

 Let $y=\dfrac{1}{x}$, $z=x+\dfrac{1}{x}.$ we have 
 \[xP(x)+\dfrac{1}{x}P\left (\dfrac{1}{x}\right )+\left ( x+\dfrac{1}{x}\right ) P\left ( x+\dfrac{1}{x}\right ) =\left (x+\dfrac{1}{x}\right )\left ( P\left (x-\dfrac{1}{x}\right )+P(-x)+P\left (\dfrac{1}{x}\right )\right ).\]

 Simplify using $P(x)=P(-x),$
 \[\left (x+\dfrac{1}{x}\right ) \left (P\left (x+\dfrac{1}{x}\right )-P\left (x-\dfrac{1}{x}\right )\right )=\dfrac{1}{x}P(x)+xP\left (\dfrac{1}{x}\right ).\]

 Expand and combine like terms, both sides are of the form 
 \[c_{n-1}x^{n-1}+c_{n-3}x^{n-3}+\cdots+c_1x+c_{-1}x^{-1}+\cdots+c_{-n+1}x^{-n+1}.\]

 They have the same values for infinitely many $x.$ They must be identical. We just compare their leading terms. On the left hand side it is $2nx^{n-1}$. There are two cases for the right hand sides: If $n>2$, it is $x^{n-1}$; If $n=2$, it is $(1+a_0)x.$ It does not work for $n>2.$ When $n=2,$ we have $4=1+a_0.$ therefore $a_0=3.$

 The solution: $P(x)=c(x^2+3)$ for any constant $c.$ 

 -JZ

 


\end{enumerate}